% Part: incompleteness
% Chapter: incompleteness-provability
% Section: godels-paper

\documentclass[../../../include/open-logic-section]{subfiles}

\begin{document}

\olfileid{inc}{inp}{gop}

\olsection{مقارنة ببحث غودل الأصلي}

تستحق ورقة غودل الصادرة سنة 1931 قراءة متأنية؛ ففي مقدمتها عرض
لأصل الأفكار التي مرّت بنا. بل يكشف التصفح السريع ترتيب حجته:
يبدأ بوصف النسق الصوري $P$، من تركيب وبديهيات وقواعد برهان،
ثم يعرّف الدوال والعلاقات العودية البدائية، ويثبت أن $x B y$
عودية بدائية وأن هذه الدوال والعلاقات تمثل في $\Th{P}$.
بعد ذلك يبرهن عدم الاكتمال على الوجه المتقدم. وفي القسم الثالث
يبيّن إمكان اختيار العبارة غير القابلة للإثبات جملة من لغة الحساب؛
ومن هناك جاءت لمّة $\beta$ التي استعملناها كذلك في معالجة
المتتاليات وإثبات تمثيل الدوال العودية في $\Th{Q}$.
ولم يفرد غودل مجموعة دنيا من البديهيات تكفي لهذا العمل، لكننا
نعلم الآن أن $\Th{Q}$ كافية. ثم ختم في القسم الرابع بمخطط
برهان مبرهنة عدم الاكتمال الثانية.

\end{document}
